\section{The integer prime-sum operator and its application}
\label{sec:primesum}
Let \(N\) be even. For \(1\leq i,j\leq N\), first put
\[
 A_N(i,j)=\frac{\Lambda(i+j)}{N}.
\]
The normalized vector inner product is
\(\langle f,g\rangle_N=N^{-1}\sum_i\overline{f_i}g_i\);
its induced operator norm agrees with the Euclidean one.
The adjacency support used at the end will exclude both
proper prime powers and diagonal entries.

\begin{lemma}[Signed integer-input transfer]
\label{lem:integertransfer}
For every fixed \(B>0\), there is a fixed \(C>0\) such that,
with \(Q=(\log N)^C\),
\[
 \left\|A_N-
    \left(\frac1N\sum_{q\leq Q}
                \frac{\mu(q)}{\phiG(q)}c_q(i+j)\right)_{ij}
       \right\|\ll_B(\log N)^{-B}.                          \tag{16}
\]
Here \(c_q(t)=\sum_{a\bmod q,(a,q)=1}\e(at/q)\)
is the Ramanujan sum. The complete signed main term is retained.
\end{lemma}
\begin{proof}
We record the argument orientation explicitly.
For distinct symbolic vectors \(f,g\), set
\(F_f(\theta)=\sum_{i=1}^N f_i\e(i\theta)\) and
\(S_M(\theta)=\sum_{r\leq M}\Lambda(r)\e(r\theta)\),
where \(M=4N+4\). Exact Fourier inversion gives
\[
 \langle f,A_Ng\rangle_N
 =N^{-2}\int_0^1 S_M(\theta)
             \overline{F_f(\theta)}F_g(-\theta)\,d\theta .
 \tag{17}
\]
Thus the left matrix factor is the conjugate of the first
argument. Parseval and Cauchy--Schwarz give
\[
 \int_0^1|F_f(\theta)F_g(-\theta)|\,d\theta
                  \leq N\|f\|_N\|g\|_N.                    \tag{18}
\]

Use major arcs \(|\theta-a/q|\leq T/N\), \(q\leq Q\),
where \(T=(\log N)^D\) and \(D\) is a sufficiently large
fixed exponent. These arcs are disjoint for large \(N\).
Take the integer \(H=\lceil N/T\rceil\).
Dirichlet approximation supplies coprime \(a,q\) with
\(1\leq q\leq H\) and
\[
 |\theta-a/q|\leq1/(qH)\leq T/(qN).
\]
If \(q\leq Q\), this lies in the stated major arc.
On their complement \(Q<q\leq H\leq2N/T\), and
Vaughan's standard estimate~\cite{Vaughan} implies
\[
 |S_M(\theta)|
 \ll N(\log N)^4
       \left(Q^{-1/2}+T^{-1/2}+N^{-1/5}\right).
\]
For completeness, the estimate being used is
\[
 |S_M(\theta)|\ll
 (Mq^{-1/2}+M^{4/5}+M^{1/2}q^{1/2})(\log M)^4
\]
when \(|\theta-a/q|\leq q^{-2}\).
The approximation also satisfies \(|\theta-a/q|\leq q^{-2}\),
because \(q\leq H\). On the minor arcs its denominator
satisfies \(Q<q\leq2N/T\).
Equation (18) therefore turns this into an operator error
smaller than \((\log N)^{-B}\) if \(C,D\) are large enough.

On each major arc, Siegel--Walfisz and partial summation give,
with arbitrary prescribed logarithmic accuracy,
\[
 S_M(a/q+\beta)
  =\frac{\mu(q)}{\phiG(q)}V_M(\beta)
          +O(N(\log N)^{-A}),\qquad
 V_M(\beta)=\int_0^M\e(t\beta)\,dt .
\]
The ranges \(q\leq Q\) and \(|\beta|\leq T/N\) are fixed
logarithmic ranges. Choose the underlying prime estimate
with accuracy sufficient to absorb the partial-summation
factor \(1+M|\beta|\).
The uniform error, inserted in (17)--(18), is again
\(O((\log N)^{-B})\).

It remains to justify the entire retained main kernel,
including the arc tails.
For \(s=i+j\in[2,2N]\), Fourier inversion of the interval
\([0,M]\), followed by one integration by parts in each
oscillatory tail, gives
\[
 \left|\int_{-T/N}^{T/N}
            V_M(\beta)\e(-s\beta)\,d\beta-1\right|
 \leq C\left(\frac{N}{T(i+j)}+\frac1T\right).
\]
The normalized matrix \(N^{-1}\) times this error has norm
\(O(1/T)\): the Hilbert matrix \((1/(i+j))\) has bounded
norm, for example by the Schur test with weights \(i^{-1/2}\).
The phase at \(a/q\) is a diagonal unitary on each side,
so does not change this bound.
Summing the explicit remainder over reduced \(a\bmod q\)
and \(q\leq Q\) costs at most \(O(Q/T)\).
Choose \(D>C+B+10\).
The retained phases sum to
\(\mu(q)c_q(-(i+j))/\phiG(q)\), which equals the displayed
main term since \(c_q(-s)=c_q(s)\).
This proves (16) without applying an unsigned estimate to
that main term.
\end{proof}

\paragraph{Formalized whole-prefix minor-grid input.}
The maintained owners \texttt{HamiltonianPrimeMinorArc} and
\texttt{HamiltonianPrimePrefixMinor} consume the existing actual
coefficient-mean Vaughan theorem and its proved physical cutoffs.
Precisely, let \(M\geq2^{18}\), \(1\leq P\leq M^{1/40}\), and
let the positive auxiliary period \(R\) be divisible by \(M\)
and by every integer \(1\leq q\leq P\). Put
\[
 W=(R/M)P,\qquad
 \mathfrak M_R(P,W)=
 \left\{\frac{Ra}{q}+j\pmod R:
  1\leq q\leq P,\ 0\leq a<q,\ (a,q)=1,\ |j|\leq W\right\}.
\]
For one absolute \(C>0\), independent of every displayed
parameter, the complete actual prefix satisfies
\[
 \left|\sum_{r=1}^M\Lambda(r)\e(kr/R)\right|^2
 \leq C\left(\frac{M^2}{P}+M^{8/5}\right)
                     (1+\log(2M))^{10},
 \qquad k\in\mathbb Z,\quad k\bmod R\notin\mathfrak M_R(P,W).
\]
The initial Vaughan cutoff is retained and bounded explicitly;
every prime power remains in the prefix. There is no divisor
power loss in this estimate. All five new declarations pass
the focused foundational-axiom check.
The full sampled-operator transfer, retained major-arc main term
and spectral comparison remain formalization obligations.

\begin{lemma}[The nontrivial ratio is one half]
\label{lem:integerspectrum}
Apart from at most one eigenvalue near \(1\) and at most one
near \(-1\), every eigenvalue of \(A_N\) has modulus at most
\(1/2+o(1)\). More precisely, in decreasing algebraic order,
\[
 \lambda_2(A_N)\leq1/2+o(1),\qquad
 \lambda_{N-1}(A_N)\geq-1/2-o(1).
 \tag{19}
\]
\end{lemma}
\begin{proof}
Use the orthonormal model basis indexed by reduced fractions
\(a/q\bmod1\), with squarefree \(q\leq Q\), including
\(0/1\) and \(1/2\).
On this basis define the self-adjoint operator
\[
 C_Q e_{a/q}=\frac{\mu(q)}{\phiG(q)}e_{-a/q}.
\]
Its eigenvalues at \(0/1,1/2\) are \(1,-1\).
All other blocks have norm \(1/\phiG(q)\leq1/2\).
It is the finite cutoff of a compact direct sum of these
orthogonal blocks.

Let \(E_Qe_{a/q}(j)=\e(aj/q)\), using the normalized sample
inner product. Its Gram matrix has diagonal 1.
For two different fractions, geometric summation bounds
the off-diagonal entry by \(C Q^2/N\).
There are \(O(Q^2)\) basis vectors, so
\[
 \|E_Q^*E_Q-I\|\ll Q^4/N=o(1).
\]
Direct multiplication, with the first argument on the left
as in (17), gives
\[
 (E_QC_QE_Q^*)(i,j)
       =\frac1N\sum_{q\leq Q}
                    \frac{\mu(q)}{\phiG(q)}c_q(i+j).
\]
The polar map \(U_Q=E_Q(E_Q^*E_Q)^{-1/2}\) is an isometry,
and the evaluated operator differs from
\(U_QC_QU_Q^*\) by \(O(Q^4/N)\) in operator norm.
Append zero eigenvalues to this finite model.
Equation (16) and Weyl's inequality give (19).
There is one exceptional eigenvalue at each end of the
spectrum, not two positive exceptions.
\end{proof}

\begin{lemma}[Uniform degrees and a true prime support]
\label{lem:primeweights}
Define the true prime-edge weight matrix
\[
 Q_N(i,j)=\frac{\log(i+j)}{N}
       \one_{\,i\ne j,\ i+j\ {\rm prime}}.
\]
It is symmetric and supported between the odd and even
integers, which form equal parts.
For every fixed \(B>0\),
\[
 \|Q_N\one-\one\|_\infty\ll_B(\log N)^{-B},\qquad
 \|Q_N-A_N\|\ll N^{-1/2}(\log N)^2,\qquad
 \max_{ij}Q_N(i,j)\leq\frac{\log(2N)}N.                      \tag{20}
\]
If \(D_N=\diag(Q_N\one)\), its normalized operator
\(D_N^{-1/2}Q_ND_N^{-1/2}\) has simple eigenvalues \(1,-1\)
and all others of modulus at most \(1/2+o(1)\).
\end{lemma}
\begin{proof}
The exact raw row sum is
\[
 (A_N\one)_i
       =\frac{\psi(i+N)-\psi(i)}{N}.
\]
The classical strong prime number theorem,
\(\psi(x)=x+O_B(x/(\log(x+2))^B)\), gives \(1+
O_B((\log N)^{-B})\) uniformly for \(1\leq i\leq N\).
Constants for bounded \(x\) can be absorbed; equivalently
split the range before applying the estimate.

For a fixed input, every proper-prime-power output
determines at most one other vertex.
There are \(O(N^{1/2}\log N)\) such outputs up to \(2N\),
each of weight at most \(\log(2N)\).
The normalized removed row sum is therefore
\(O(N^{-1/2}(\log N)^2)\).
The diagonal costs at most \(\log(2N)/N\) per row.
Symmetry and the row-sum norm bound prove the second
estimate in (20); the same errors preserve the first.
The entry bound is immediate.
Every off-diagonal prime output is odd, because the sole
representation of prime 2 by positive equal-bounded inputs
is the removed loop at vertex 1.
Thus the remaining support is exactly bipartite.

Its degree matrix is \(I+o(1)\), so normalization changes
the operator by \(o(1)\).
The normalized matrix has eigenvalues \(1,-1\), with vectors
\(D_N^{1/2}\one,D_N^{1/2}s\).
Equation (19) permits no further eigenvalue of modulus
larger than \(1/2+o(1)\).
It also forces these two eigenvalues to be simple;
multiple components would supply extra eigenvalues at 1.
This proves the final assertion.
\end{proof}

\begin{proof}[Proof of Theorem~\ref{thm:filz}]
Apply Lemma~\ref{lem:scaling} to \(Q_N\), with any fixed gap
smaller than \(1/2\), using (20).
For instance take the row accuracy \(B=10\);
then \(\varepsilon L^2\to0\).
The result is a symmetric stochastic bipartite \(P_N\)
with exactly the true prime-edge support and
\[
 \|P_N-Q_N\|=o(1),\qquad
 \max_{ij}(P_N)_{ij}\ll\frac{\log N}{N}.
\]
Its nontrivial spectral ratio is \(1/2+o(1)\).
Fix a positive gap such as \(1/3\) for all sufficiently
large \(N\), and take \(d_N=\lfloor N^{1/4}\rfloor\).
The finite rounding theorem applies, since
\[
 d_N\alpha_N L^2\ll N^{-3/4}(\log N)^3\longrightarrow0.
\]
Its complete operator error is
\[
 O\left(N^{-3/4}(\log N)^2+
             N^{-1/8}\sqrt{\log N}+N^{-1/4}\log N\right)=o(1).
\]
It gives a simple \(d_N\)-regular spanning subgraph of
\(G_N\), whose second eigenvalue is at most
\((1/2+o(1))d_N\).
This is (4).
For sufficiently large \(N\) its gap is at least \(1/3\).
The polynomial degree exceeds \((9\log N)^K\).
The established bipartite theorem (2) therefore supplies
a Hamilton cycle in this subgraph, and hence in \(G_N\).
All proper-prime-power and loop contributions were removed
before the support was used.
\end{proof}

\begin{remark}[Quantifiers and the finite remainder]
Every sufficiently large even \(N\) is covered, rather than
a selected subsequence of prime gaps.
The starting threshold depends on classical prime estimates
and on the large-order threshold in (2), and has not been
made numerical. This draft therefore leaves the finite
remainder of the all-orders Filz conjecture unresolved.
No finite prime experiment is used to fill that remainder.
\end{remark}
